\documentclass[11pt] {report}
\usepackage[onehalfspacing]{setspace}
\usepackage{arial}
\usepackage[margin=2.5cm]{geometry}
\usepackage{graphicx} %for image inclusion
\usepackage{natbib} %for organising references 
\usepackage{hyperref} %for making references clickable
\usepackage{amsmath} % for mathematical notations

\begin{document}

%cover page
\title{Initial Thesis Plan: How does Cyclical Change Impact Functional Diversity, Ecosystem Response and Resilience across Ecoregions?} 
\author{Joe knox}
\date{\ 26/10/2026}

\maketitle

\pagenumbering{⟨arabic⟩}

%insert table of contents
\newpage
\tableofcontents

%begin report
\newpage
\chapter{Abstract}
\label{chap:abstract}


\begin{itemize}
 
\item 300 words
\item set scene and introduce topic
\item explain research, working hypothesis and main knowledge gap
\item summarize work programme (methods) without extensive details

\end{itemize}
 

\noindent


\chapter{Introduction}
\label{chap:introduction}

\section{Background}

\noindent

\begin{itemize}

\item 3000 words max
\item what do we know?, set work in wider global context
\item critical analysis of exisiting literature
\item what do we not know? (knowledge gap and novelty of your work)
\item allow reader to understand scope of project
\item references and ~2-3 figures/schematics/diagrames
\item divided into subsections

\item Biodiversity is not static 
 \item Species composition with communities across seasons, with migration and hibernation, and even from day to night (daily patterns) as they go to sleep
 \item A lot of work has been done looking shift in biodiversity across space e.g.latitudinal shifts, and along environmental gradients but not so much across the annual cycle, particularly outside of coarse seasonal change (i.e. winter vs summer). This 
 \item By looking at the contributions of species to functional diversity in different seasons, and acitivty periods we can assess how  they might be contributing to ecosystem services (ie. by combining with diet types and using traits as a proxy for ecoystem functioning)
 \item But at what granularity of cyclical change can we see functional shifts and how fine-scale do we go in terms of defining cyclical patterns? daily, weekly, monthly and does this matter for global change?
 \item - important for characterising niche occupancy shifts from day to night and species interactions i.e. resource competition, niche partioning (limiting similarity hypothesis)
\item novelty- nobody has focused on cyclical change in such detail, outside of monthly turnover or single species daily acitivty pattens. for birds acitivty periods is quite novel but not so much seasonality, for mammals hibernation an migration is quite novel, whilst actiivty period is more well-researched?. 
\item For other taxa, seasonal functional turnover studies are essentially non-existent

\end{itemize}

\chapter{Aims and Objectives}
\label{chap:aims and objectives}

\begin{itemize}

\item 300 words

\end{itemize}

\section{Aims}

\noindent

\begin{itemize}
\item 1-3 aims, that address knowledge gaps previously stated. one sentence per hypothesis
\item chapter 1: assess the impacts of seasonal/monthly cyclical changes on avian functional diversity (migration)
\item chapter 2: assess the impacts of daily cyclical change on avian functional diversity (activity period)
\item chapter 3: assess the application of characterising cyclical functional changes for other non-avian taxa mammalian/reptilian/amphibian (migration/hibernation and activity period)
\item chapter 4: understand how climate change impacts functional turnover due to cyclical changes in  birds/mammals/chosen taxa
\end{itemize}

\section{Hypotheses}

 \textbf{Hypothesis 1: In birds, migration throughout the annual cycle is important for ecosystem functioning. However, migrant functional contributions will vary with season, dietary/foraging niche, duration of stay, ecoregion and latitude} 

 %JK to TG: will seperate into seperate hypotheses for the actual paper but limited by initial thesis plan specifications, or maybe better to frame as a question rather than a hypothesis i.e. what is the impact of migration on avian fucntional diversity?

 %JK to TG: note that fallback to non-refined range for passage species (as species' AOH maps do not account for passage ranges) may result in a bias towards passage species in assemblages (which may overestimate their contribution to traitspace)

 %JK to TG: have included reasoning FYI below but i dont think i need to include this in the thesis plan

 \begin{itemize}
 \item on the one hand may expect migrants to be contributing specialist functions since they tend to specialise on particualar resources which results in the evolution of migration as a strategy in the first place, or may have particualr climatic niches which neccessitate migration across seasons/wet/dry periods
 \item migration also requires unique traits e.g. HWI/bodymass (lots of evidence for this; but see Williams and Curley, 2025)
 \item also migrants need to be somewhat divergent in their traits - limiting similarity hypothesis, to avoid competition with residents/other migrants?
 \item may also expect migrants to be more spread out, as they are likely less generalist and plastic (residents may switch their niche across seasons), occupying more extreme trait combinations
 
 \item However, there is also likely to be a lot of overlap with resident communities as both migrants and residents are captilising on similar resource peaks that are more abundant at particular times of year. e.g. an african bird migrating north for seeds will not suddenly eat fish, it will probably eat the abundant seeds that the residents are likely also captilising on. Monopolisation of these resources is a different quesiotn e.g. Williams and Curley, 2025
\item This functional overlap may also be the case across seasons- evidence for biotic homogenisation of migrants across seasons, probably due to shifting migrant species distributions and changing resource peaks (i.e. phenological mismatch).
 
\item Additionally, more transient migrants are likely to be more unique as they may be migrating long distances resulting in the need for short stop-overs (calculate migratory distance as well to confirm this). Thus, these short-stay species are likely have diet/foraging types included in their niche that are not present in their stop-over region, and therefore they may have unique trait combinations that are distinct from resident communities/long-stay migrants that are not utilising these "foreign resources" in other locations/parts of the year 
 \item however, this is unlikely to be significant for ecosystem functioning because a species present in a month is unlikely to make a signifcant functional contribution

 \item functional contribution of migrants will differ across seasons, dietary/foraging types, and ecoregions
 \item i.e. previously granivores/frugivores found to dominate in winter, whilst insectivores are more prevalent in summer (Keyser et al., 2024; Bellve et al., 2025)
\item additionally, nestedness vs overlap with resident communities tends to differ across habitat types (Oliveira et al., 2023)

\item there is also a latitudinal effect of seasonality, i.e. seasonal and resource variation is greater at northern latitudes and migration is more dominant in the northern hemisphere,
\item thus migrants are likely to contribute more to functional diversity at higher latitudes, potentially introducing more specialist functions.
\item On the other hand competition may be greater in the tropics where species richness is greatest, therefore migrants and residents may be able to better partition resources in seasonal environments (higher latitudes) than non-seasonal environments.
\item Thus, migrants may be more specialist in the tropics due to the limiting similarity hypothesis (although greater producitivty may also facilitate niche packing of migrants and residents)
\end{itemize}

\textbf{Hypothesis 2: In birds, ecosystem functioning will shift across the daily cycle and birds unique to the nocturnal or diurnal niche may represent unique functions.}

\begin{itemize}
\item species that are purely nocturnal/diurnal likely contribute unique and/or specialist traits e.g. beak shape for nectar (lots of diurnal species) vs rodents (lots of nocturnal species) will differ and species will be using different senses e.g. heat vs colour vision which will impact their trait shape/size and this will impact ecosystem functioning, niche space and species interactions.
\item species existing in both night and day likely feed on non-colour reliant resources. or may be more plastic in their niche (generalists)
\end{itemize}

\textbf{Hypothesis 3: In mammals/other taxa, changes across the annual cycle (i.e. migration) will likely play as less significant role than in birds (although hibernation in mammals may be important), however daily functional shifts may be more important.}

\begin{itemize}
\item migration is much more widespread in birds (~40% of birds migrate whilst only ~20%  mammals migrate- ungulates and marine mammals are really the only migrating mammals and in amphibians, migration is will be more specific to breeding season and probably much more relevant for local aquatic functional diveristy (ponds))
\item but daily changes are likely more important as activity period seems to be more important part of ecology in other taxa, with more nocturnal mammals? (see Cox and King et al papers)
\end{itemize}

\textbf{Hypothesis 4: Climate-induced shifts in species' distributions will disproportionately effect migrants and species specific to one acitvity period as they may be more specialist and less adaptable under climate change, impacting ecosystem functioning}

\item migrants/activity period specialists likely have more specific climate-niches than residents, driving their migratory/acitvity period life-history in the first place. Thus, they may be more sensitive than residents or acitvity period generalists to climate change.
\item if they are found to contribute unique functions (in above chapters), this may have real consequences for ecosystems. 
\item Additionally, the migrant/ acitivty period species gained in ecoregions are likely to be the more plastic niche switchers/ generalists which will probably be more functionally similar to the existing resident community/species existing across both nocturnal and diurnal niches



\section{Research Objectives}

\begin{itemize}

\item 3-6 objectives, one sentence per objective., Main research goals you need to address your questions/hypotheses
\item research goals you need to address your hypotheses

\end{itemize}

\textbf{chapter 1:}

\noindent

\textbf{global migratory status analysis}

\begin{enumerate}
\item  What is the contribution of migrants vs residents to annual functional diversity? migrants vs residents in local annual trait space (ecoregion)
\item  What is the contribution of migrants vs residents in each season (winter vs summer vs spring and autumn for Britain and Ireland i.e. seasonal local trait space (ecoregion))

\end{enumerate}

\textbf{ migratory status as a gradient analysis}

This analysis is necessary because the amount of time that a species is present in an ecoregion is not the same across migrants. Many are passage, with just a small presence and so their functional contribution may be small, but it could be unique. This may not be a problem since a migrant present for only a month/week may not make a significant contribution to the ecosystem, but it could be?
%>% n.b. note limitation with passage range fulling back to unrefined map- assemblage presence could be biased towards passage species

\begin{enumerate}
\item  What is the contribution of transient vs long-stay migrants and residents in annual ecoregional trait space? migratory status grouped by number of months present (residents are present month=12) to local annual trait space (ecoregion)
\item  What is the contribution of transient vs long-stay migrants and residents in seasonal ecoregional (winter vs summer vs spring vs autmnm for Britain and Ireland i.e. seasonal local trait space? 
\item Could extend to weekly/daily (ebird data at eco-regional scale or Zenodo for Britain and Ireland at non-ecoregional scale but not extendable globally)
\end{enumerate}


\textbf{migratory status by diet/foraging type analysis (Avoniche)}

\textbf{migratory status by ecoregion analysis (analysing across regions- are the results consistent?)}

\textbf{migratory status by contintent/latitude analysis (extending beyong Britain and Ireland, do the results still stand?)}


\textbf{Initial research plan} 


\begin{itemize}
\item 800 words
\item Outline of research intent for chapter 1 (8-15 months of research or maybe less)
\item state which hypothesis you are addressing
\item describe analysis/modelling approach and methods
\item how you intend to interpret the results/ draw conclusions 
\item diagrams/schematics useful with legends
\end{itemize}

\item addressing hypothesis 1  %note across region extension, can update later if needed
\end{itemize}

\textbf{Chapter 1: The effect of Seasonal/monthly/weekly cyclical change on functional diversity starting with birds because that is where the data is.}

\begin{itemize}

\item \textbf{Seasonal/monthly/weekly cyclical change:} migrants vs residents across the annual cycle (probably seasonally to weekly to monthly) with transience defined i.e. migratory status as a scale from presence in one day/365- 364/365 days i.e. do passage migrants contribute unique functions in the week they are present?
 \item focus on functional uniqueness: Are migrants functionally unique compared to residents
 \item what functions do they contribute?
 \item increasing granularity, are more migratory species (i.e. more transient) more functionally unique
 \item functional uniquness important as more unique, specialist migrant species may be more at risk from phenological shifts under things like climate change or more vulnerable to land-use change than generalist residents due to phenological shifts as their specialist diet types are converveted for example, resulting in phenological mismatch. So they may be lost more easily from ecosystems and their keystonefunctions may be more liekly to be lost. 
 \item Not only this but they may be more irreplaceable than specialist residents as they may provide keystone functions that are more transient (present fewer weeks of the year), so more likely to be lost.
 \item migrants found to be paritcularly important for ecosystem functioning in highly seasonal environments e.g. (keyser et al., 2024) as most migration happens in northern hemisphere. And this concerning as they are also particularly vulnerable to climate change etc... (cite)
 \item But what about other regions? boreal vs sahara for example or tropical ecoregions?

\end{itemize}



\chapter{References}
\label{chap:references}


\end{document}

